\documentclass[10pt]{article}

\usepackage[margin=0.75in,top=0.7in,bottom=0.7in]{geometry}
\usepackage[T1]{fontenc}
\usepackage{lmodern}
\usepackage{microtype}
\usepackage{booktabs}
\usepackage{array}
\usepackage{enumitem}
\usepackage{xcolor}
\usepackage[hidelinks]{hyperref}
\usepackage{titlesec}

\setlist{nosep, leftmargin=1.2em}
\setlength{\parskip}{0.35em}
\setlength{\parindent}{0pt}
\setlength{\tabcolsep}{4pt}
\frenchspacing
\emergencystretch=3em
\titleformat{\section}{\normalsize\bfseries}{\thesection.}{0.4em}{}
\titlespacing*{\section}{0pt}{0.9em}{0.2em}
\titleformat{\subsection}{\small\bfseries}{\thesubsection}{0.4em}{}
\titlespacing*{\subsection}{0pt}{0.6em}{0.1em}

% Everything in red depends on the final run (#15, Nov 13) or is still to be written.
\newcommand{\todo}[1]{\textcolor{red}{[#1]}}

\begin{document}

\begin{center}
{\fontsize{13}{16}\selectfont\bfseries VAKRA-Advanced: Tool-Using LLM Agents for Multi-Hop, Multi-Source Reasoning}\\[0.4em]
{\normalsize Final Report}\\[0.4em]
{\small Anurag Kaushal (2025202013) \quad Aayush Pandey (2025201058) \quad Piyush Agarwal (2025201082)}
\end{center}

\section{Introduction}
VAKRA-Advanced (SemEval-2027 Task 11) asks an agent to answer questions by calling tools in a
live environment: 8,000+ APIs over real databases in 62 domains, plus document collections.
There are four tasks: (1) chaining generic data operations, (2) selecting the right endpoint,
(3) multi-hop API reasoning, and (4) multi-turn dialogues that mix APIs and documents under
tool-use policies. The official baseline is a LangGraph ReAct agent that shortlists up to 288
tools per domain and lets the LLM call them in a loop.

Our Phase~1 report measured that baseline and diagnosed its failures. This report describes what
we changed and what each change gained. In short:
\begin{itemize}
\item \textbf{Three bugs in the official runner} cost more than any modelling choice. A tool-schema
  bug rejected correct calls on 22.5\% of Task~2 items, a stale tool list made Task~1 unsolvable
  for most items, and long dialogue histories overflowed the context on 17\% of Task~4 dialogues.
  We fixed all three.
\item \textbf{Tool retrieval is solved at the official shortlist size.} A hybrid BM25 + dense
  shortlist keeps every gold tool of every solvable item at $k{=}128$, so retrieval is not where
  items are lost. Its advantage is at small $k$.
\item \textbf{Four agent-level changes} target the failures that remain: rewriting the final answer
  from the tool evidence only (V5), following Task~4's tool policies by construction (V6), a
  system rule to try the tools before answering (V7), and hints for case-sensitive values (V8).
\item On our evaluation set, the final system scores \todo{final scores} on Tasks~1--4
  (macro \todo{final macro}), against 0.408, 0.343 and 0.350 for the baseline on Tasks~2--4; the
  baseline cannot complete Task~1 items at all. The largest remaining failure is choosing the wrong tool from a correct shortlist:
  \todo{47--67\%} of the failed items in Tasks~2--4.
\end{itemize}

\section{Experimental setup}
\textbf{Data.} We hold out 9 training domains from \emph{every} task (\texttt{authors},
\texttt{cars}, \texttt{european\_football\_1}, \texttt{hockey}, \texttt{ice\_hockey\_draft},
\texttt{menu}, \texttt{sales\_in\_weather}, \texttt{student\_loan}, \texttt{world}), since Tasks~2--4 share each domain's
tools and a per-task hold-out would leak them. All scores are on \texttt{dev-eval}, a fixed,
domain-stratified sample of these held-out domains: 60, 120, 35 and 60 items for Tasks~1--4.
Access to the official test set was requested and is still pending, so every number here is on
held-out \emph{training} domains.

\textbf{Model and scoring.} The agent and both judges use \texttt{gpt-oss-120b} on Groq at
temperature~0. The official evaluator replays the agent's tool calls and scores the final turn
0/1: the agent must reproduce the gold tool responses or convince a correctness judge, and a
second judge must find the answer grounded in the tool outputs. When the gold answer is a
refusal, a refusal with no tool call scores. Task~4 is scored on the final turn of each
dialogue; earlier turns are given as history with their gold answers.

\textbf{Noise.} With these sample sizes the 95\% half-widths are $\pm$0.13, $\pm$0.09,
$\pm$0.17 and $\pm$0.12 for Tasks~1--4. Temperature~0 does not make runs repeatable: between two
runs with identical input, about 10\% of items flip. A difference of a few items is therefore
not evidence, and we judge each change item by item on the items it targets, not on totals.

\textbf{Cost.} All runs together, judge included, cost \todo{\$9.36 + final run} on Groq's paid
tier.

\section{System}
We keep the official agent, runner and evaluator unchanged. Every change is a run-time patch,
applied by our script \texttt{run\_baseline.py} before the official runner starts. Each patch is
behind its own flag, and the run id records which flags were on. This keeps the official baseline
reproducible and makes every ablation a one-flag difference. We dropped our plan to write our
own agent (V1): the patches already carry every fix, and a separate agent would have needed a
parity run we could not afford. The cost is that changes needing control over the agent loop
itself are out of scope.

\begin{center}\small
\begin{tabular}{@{}lllp{0.52\textwidth}@{}}
\toprule
Change & Flag & Tasks & What it does \\
\midrule
\multicolumn{4}{@{}l}{\emph{Fixes to the official runner}} \\
Schema fix & \texttt{-{}-fix-tool-schema} & all & Tools without parameters get an empty schema
  instead of one with a required internal field, which Groq rejected \\
Tool refresh & \texttt{-{}-refresh-tools} & 1 & Re-lists the tools after Task~1's per-item tool
  set switches \\
Short history & \texttt{-{}-short-history} & 4 & Earlier answers over 2,000 characters keep their
  first rows plus a note with the total \\
\midrule
\multicolumn{4}{@{}l}{\emph{Variants}} \\
V2 & \texttt{-{}-retriever hybrid} & 2--4 & BM25 + dense tool shortlist (reciprocal-rank fusion) \\
V3 & \texttt{-{}-repair-tool-calls} & all & A tool call the API rejects is retried once, with
  the error shown to the model \\
V4 & -- & 3 & Per-hop tool retrieval; dropped after an offline upper bound showed no gain \\
V5 & \texttt{-{}-evidence-answer} & all & The final answer is rewritten from the tool results
  only \\
V6 & \texttt{-{}-policy-guard} & 4 & Task~4 tool policies are enforced, not only stated \\
V7 & \texttt{-{}-tools-first} & 2--3 & A system rule: try the tools before answering, never ask
  the user back \\
V8 & \texttt{-{}-tool-hints} & 2--3 & V7 plus hints on case-sensitive values and multi-hop
  questions \\
\bottomrule
\end{tabular}\\[0.2em]
{\footnotesize \textbf{Table 1.} Our changes. All are patches applied at run time; the official
code is unchanged.}
\end{center}

\section{Results}
\textbf{Final configuration.} \todo{Table 2: one row for the official baseline (V0) and one for
the final run with all flags (Table 1 except V4), per task with 95\% half-widths and the macro
average; run ids in Section~9. Fill in from the final run on Nov 13.}

\textbf{Ablation.} Table~3 adds one change at a time. Each row is its own run on
\texttt{dev-eval}; a dash means the change does not apply to that task and the score carries
over from the row above. V3 is scored only as part of the final run.

\begin{center}\small
\begin{tabular}{@{}lccccc@{}}
\toprule
 & Task 1 & Task 2 & Task 3 & Task 4 & Macro \\
 & ($n{=}60$) & ($n{=}120$) & ($n{=}35$) & ($n{=}60$) & \\
\midrule
V0, official baseline                 & --$^{\dagger}$ & 0.408 & 0.343 & 0.350$^{\ast}$ & -- \\
\quad + tool refresh                  & 0.483 & -- & -- & -- & -- \\
V2 hybrid shortlist ($k{=}128$)       & -- & 0.450 & 0.371 & -- & -- \\
\quad + schema fix + tool refresh     & 0.450 & 0.558 & 0.486 & 0.317$^{\ast}$ & 0.453 \\
\quad + V5 evidence-only answer       & 0.583 & 0.583 & 0.486 & 0.317$^{\ast}$ & 0.492 \\
\quad + short history                 & -- & -- & -- & 0.433$^{\ast}$ & 0.521 \\
\quad + V6 policy guard               & -- & -- & -- & 0.500 & 0.538 \\
\quad + V7 tools-first rule           & -- & 0.633 & 0.486 & -- & 0.551 \\
\quad + V8 hints                      & -- & \textbf{0.650} & 0.486 & -- & \textbf{0.555} \\
\midrule
95\% half-width                       & $\pm$0.13 & $\pm$0.09 & $\pm$0.17 & $\pm$0.12 & \\
\bottomrule
\end{tabular}\\[0.2em]
{\footnotesize \textbf{Table 3.} One change at a time on \texttt{dev-eval}.
$^{\dagger}$Stopped after 7 items, all of which hit the iteration limit because of the stale
tool list (Section~6). $^{\ast}$Includes dialogues credited after the agent crashed (gold is a
refusal and no tool was called): 3 for V0, 1 for the next rows. V6 refuses on purpose instead.}
\end{center}

Most single steps in Table~3 are smaller than their task's 95\% half-width, so we do not rely
on totals. For each change, Sections~5--7 report what happened
to the items it targets.

\section{Tool retrieval}
The baseline gives the LLM only the top $k$ of a domain's tools ($k{=}128$ by default), ranked
by MiniLM embeddings of the question against each tool's \texttt{name: description}. We treat
this as ad-hoc retrieval and count an item at $k$ only if \emph{all} its gold tools are in the
top $k$, so recall@$k$ caps the agent's score. We compared the dense retriever with BM25
($k_1{=}1.5$, $b{=}0.75$) and their reciprocal-rank fusion (constant 60).

\begin{center}\small
\begin{tabular}{@{}llcccccc@{}}
\toprule
Task & Retriever & R@10 & R@20 & R@32 & R@50 & R@128 & MRR \\
\midrule
2 & Dense (MiniLM, baseline) & \textbf{.825} & .900 & .925 & .958 & \textbf{.992} & .563 \\
  & BM25                     & .750 & .842 & .933 & .950 & \textbf{.992} & .477 \\
  & Hybrid (V2)              & .817 & \textbf{.925} & \textbf{.958} & \textbf{.975} & \textbf{.992} & \textbf{.565} \\
\midrule
3 & Dense (MiniLM, baseline) & .343 & .514 & .714 & .800 & .943 & .624 \\
  & BM25                     & \textbf{.371} & .543 & .714 & \textbf{.971} & \textbf{1.00} & .540 \\
  & Hybrid (V2)              & .314 & \textbf{.571} & \textbf{.800} & .857 & \textbf{1.00} & \textbf{.685} \\
\bottomrule
\end{tabular}\\[0.2em]
{\footnotesize \textbf{Table 4.} Gold-tool recall of the shortlist on \texttt{dev-eval}
(120 / 35 items). No LLM calls.}
\end{center}

\begin{itemize}
\item \textbf{At $k{=}128$, retrieval is solved.} The hybrid shortlist keeps every gold tool of
  every solvable item. The one Task~2 miss is an item whose gold tool is not in the domain's
  catalog, so no agent can solve it. None of the items lost in our runs is a retrieval failure.
\item \textbf{The hybrid's advantage is at small $k$.} At $k{=}32$ it keeps the gold tools for
  95.8\% of Task~2 and 80.0\% of Task~3 items, against 92.5\% and 71.4\% for the baseline's
  retriever. At $k{=}50$, BM25 alone is best on Task~3 (97.1\%): multi-hop questions name later
  hops by keyword more than by meaning.
\item \textbf{Smaller shortlists still cost accuracy end to end.} Before the runner fixes,
  going from 128 to 50 hybrid tools lost 3 Task~2 and 3 Task~3 items, in line with Task~3's
  recall drop. We therefore keep $k{=}128$, where V2's end-to-end gain over the baseline
  (+5 Task~2 items, +1 Task~3) is within noise.
\item \textbf{Per-hop retrieval (V4) has nothing to gain.} We re-ranked the tools at each hop
  using the question plus the previous hop's \emph{gold} call and response, an upper bound for a
  live agent. It changes at most 1 of the 35 Task~3 items at any $k$. Of the 7 items that miss a
  gold tool at $k{=}32$, 4 miss it at hop~1, which per-hop retrieval cannot change, and the
  previous responses are mostly bare values (\texttt{[["USA"]]}) that say little about the next
  tool. We did not build a live V4.
\end{itemize}

\section{Tasks 1--3: single-turn questions}

\subsection*{Runner fixes}
\textbf{Tool schema.} In Phase~1, 22.5\% of Task~2 and 17.1\% of Task~3 items ended with the API
rejecting a tool call. For tools without parameters the official wrapper passes no schema, so
LangChain infers one with a required internal \texttt{config} field, and Groq's strict
validation rejects the model's correct call. With the fix, 18 of the 25 Task~2 items the bug had
killed and all 3 Task~3 items are correct, and rejected calls fall to 1.7\% and 5.7\% of items.
What remains are genuine model mistakes, such as a missing required argument, which V3 targets.

\textbf{Stale tool list (Task~1).} Task~1's server switches between two tool sets when an item
loads its data: generic operations (\texttt{filter\_data}, \texttt{sort\_data}, \ldots) or
item-specific ones (\texttt{select\_data\_equal\_to}, \ldots). The official runner lists the
tools once per domain and never again, so on items of the other kind every call returns
``Unknown tool'' and the agent loops until a limit. This affects 48 of our 60 Task~1 items, and
600 of the 2,085 items in the official test mapping. Re-listing the tools before each item
removes every ``Unknown tool'' call and makes Task~1 scorable at all (0.483).

\subsection*{V5: answers grounded in the evidence}
After the fixes, 12.7\% of all items (35 of 275) had the right answer or the right tool responses
and failed only the groundedness judge. Reading all 35 by hand gave two main causes. In 13, the
model adds something the tools never returned: units (``200\,cc''), interpretations of codes,
world knowledge (``the Ford Torino is American''), or framing such as ``the only''. In 12, the
answer restates a filter that lives in the tool call's arguments (``these are the US models''),
which the judge cannot see: it is given the query and the tool responses, not the calls. The
second cause is a limit of the evaluator and affects every system.

V5 adds one LLM call that rewrites the final answer using only values present in the tool
results, copied verbatim, and says the information is not available only when no tool result
supports the draft. Offline, on the 35 items plus 35 correct control items, it recovered 10 of
the 25 ungrounded items in Tasks~2--4, none of which pass when their original answers are simply
re-scored, with no net loss on the controls. In the full run Task~2 rose from 0.558 to 0.583 and
Task~1 from 0.450 to 0.583. Part of Task~1's gain is likely judge noise: 5 of its 10 ungrounded
items pass on a re-score alone.

\subsection*{V7 and V8: making the agent use its tools}
After V5, 18 of 120 Task~2 and 10 of 35 Task~3 items were answered without any tool call, and
none of them correctly. The model decided that no tool fits, and in 13 of the 28 cases asked the
user to clarify, although the benchmark has no user. The official runner gives the agent no
system message in Tasks~2--3. V7 adds one, passed as LangGraph's \texttt{prompt} so that it is
not recorded in the trajectory: try the tools before answering, related tools too; never ask
back, take the most reasonable reading; say an answer is not available only when the tool
results do not contain it. Items with no tool call fall from 28 to 3 and ask-backs from 13 to 0.
Task~2 rises from 0.583 to 0.633 and Task~3 is unchanged. Most of the former no-tool items still
fail, now by calling the wrong tool: the rule removes a refusal but does not find the right tool.

The error analysis (Section~8) then showed two further patterns. In 5 items the agent called the
right tool with the right value in the wrong case (\texttt{venue='steamship'}, gold
\texttt{'STEAMSHIP'}); the data is case-sensitive, the lookup returned 0, and the agent reported
that. In 7 of 12 Task~3 wrong-tool failures, the chain stopped after the first hop. V8 adds two
sentences to V7's rule: retry an empty or zero lookup in other capitalizations, and keep calling
tools until every part of the question is answered. It fixed the 3 Task~3 items where it applied:
the agent retried \texttt{"ford torino"} in lowercase, found the country and finished the chain.
It fixed none of the 4 Task~2 items, where the agent either accepted the 0 or retried other
spellings in the same case. Task~2 rose from 0.633 to 0.650, within noise.

\textbf{V3: repairing rejected tool calls.} After the schema fix, 6 items in the V5 runs still end
with the API rejecting a call (invalid JSON or a missing required argument). V3 catches the
rejection at the LLM call, shows the model the error and its rejected call, and retries once; the
note is not kept in the history. It was tested offline against a LangGraph agent with a fake LLM
and is scored only in the final run, where \todo{n of the rejected-call crashes disappear}.

\section{Task 4: multi-turn dialogues under tool policies}

\textbf{Long histories.} The official runner puts every earlier turn's gold answer into the
history in full. In 10 of the 60 dialogues an earlier answer is a whole table, 1.3M to 9.4M
characters, and the first LLM call is rejected before any tool is called. Our fix keeps each
earlier answer over 2,000 characters to its first rows plus a note such as ``(61223 rows in
total; only the first 75 are shown)''. The largest history shrinks from about 2.3M tokens to
1.6k, and queries, turn order and policies are unchanged. Live, 6 of the 10 dialogues become
correct and the other 50 change by +1 net, i.e.\ noise: Task~4 rises from 0.317 to 0.433.

\textbf{Tool policies (V6).} Each dialogue may carry a policy restricting the use of document
retrievers (the \texttt{query\_*} tools). The official agent only states it to the model, and
for single-turn dialogues not at all; in the V5 run it followed none of them. In the 898
training dialogues every policy is one of 4 templates over 12 topics, and the gold answer is the
refusal ``I can not answer.'' \emph{exactly} when only retrievers are allowed and the policy
applies: in 65 of 65 global policies and 68 of 68 topic policies naming the domain's topic, and
in none of the 79 naming another topic. V6 enforces this: when only retrievers are allowed and
the policy applies, it refuses with no tool call; when retrievers are forbidden, the agent runs
without them. For domains outside our training split, the domain's topic comes from one LLM
call per domain, which agrees with the benchmark on 25 of 29 labelled training domains. Live,
V6 refused exactly the 11 dialogues whose gold is a refusal and nothing else, against 0
deliberate refusals by the baseline. Task~4 rises from 0.433 to 0.500; the 4 losses are on
dialogues the guard did not touch. On the official test set the guard should matter more:
IBM's policy judge, which is not public, scores a turn 0 when a policy is broken.

\section{Error analysis}
\todo{Replace every number in this section with the final run's (scripts/error\_analysis.py
-{}-run-id <final run>), and pick the examples from the annotated sample of about 20 failed items
per task. The text below is the preliminary version on the latest run of each task.}

We give each failed item one category. The first split is the evaluator's own exact-match
verdict: if every gold tool response is among the agent's, the tool part was right and the answer
failed; otherwise the reason comes from comparing the agent's calls with the gold calls.

\begin{center}\small
\begin{tabular}{@{}lcccc@{}}
\toprule
Category & Task 1 & Task 2 & Task 3 & Task 4 \\
\midrule
Agent crash                                         & 4 & 1 & 1 & 2 \\
Right answer, not grounded                          & 3 & 0 & 2 & 2 \\
Right tool results, answer not grounded             & 1 & 9 & 1 & 6 \\
Tool results missing: no tool call                  & 0 & 1 & 0 & 1 \\
Tool results missing: different tool                & 10 & \textbf{23} & \textbf{12} & \textbf{14} \\
Tool results missing: argument differs only in case & 0 & 4 & 1 & 0 \\
Tool results missing: other argument value          & 7 & 6 & 1 & 5 \\
\midrule
Failed / items                                      & 25 / 60 & 44 / 120 & 18 / 35 & 30 / 60 \\
\bottomrule
\end{tabular}\\[0.2em]
{\footnotesize \textbf{Table 5.} Failed items by category \todo{preliminary: V7 run for Tasks
2--3, V6 run for Task 4, V5 run for Task 1}. No failed item is ``right tool results, wrong
answer''.}
\end{center}

\begin{itemize}
\item \textbf{Choosing the tool is now the main failure} (47--67\% of failures in Tasks~2--4),
  although the gold tool is in the shortlist. In Task~2 the model picks a plausible neighbour
  (``shoutouts in 1977'' $\rightarrow$ \texttt{get\_sum\_sho\_by\_year} instead of a count of
  shutouts); some gold choices are hard to infer from names at all (``the car with the best crash
  protection'' $\rightarrow$ \texttt{get\_heaviest\_car\_details}). In Task~3 the chain breaks at
  the second hop. In Task~4 the agent splits one composite gold tool into many simple calls, or
  queries documents where the gold uses a data tool. For Task~1, ``different tool'' is a weak
  signal: it has many equivalent tool paths, and 26 of its 35 correct items also skip a gold tool.
\item \textbf{Argument values in the wrong case} (5 items in Tasks~2--3). The model cannot know
  the stored case from the question; V8 recovers it only when the model retries.
\item \textbf{Groundedness} is still 21\% of failures (24 of 117). Many are the evaluator limit
  from Section~6, where the answer restates a condition that lives in the tool's arguments; the
  others add a unit or an inference the tool output does not contain.
\item \textbf{Crashes are rare} (8 of 117): rejected tool calls (3, V3's target), tool responses
  too large for the API (2), and single context-length and recursion-limit failures.
\end{itemize}
\todo{Add 2--3 annotated example traces from the final sample.}

\section{Limitations and future work}
\begin{itemize}
\item \textbf{Small samples and run-to-run noise.} With 35--120 items per task and about 10\% of
  items flipping between identical runs, most single changes are within their 95\% intervals.
  Our evidence for each change is item-level: what happened to the items it targets.
\item \textbf{Held-out training domains, not the test set.} Access to the official test set is
  still pending. Our 9 held-out domains mimic unseen domains, but the test set also has newly
  curated items.
\item \textbf{One model.} Every run uses \texttt{gpt-oss-120b} on Groq, also as the judge. We did
  not test whether the changes transfer to other models.
\item \textbf{The evaluator's judges.} The groundedness judge cannot see tool arguments, which caps
  every system's score, and Task~1's groundedness verdicts are unstable on re-scoring. Task~4's
  policy judge is not public, so policy compliance is only measured indirectly.
\item \textbf{No control over the agent loop.} Because every change is a patch on the official
  agent, ideas that need a different loop, such as planning or verifying a chain before
  answering, were out of scope.
\item \textbf{Future work.} The largest failure, choosing the wrong tool from a correct
  shortlist, needs a stronger selection step: re-ranking the shortlist with the LLM, or few-shot
  examples of similar questions and their gold tools from the training split. Case-sensitive
  lookups are better fixed in the tool layer than by prompting. Tool responses too large for the
  API could be truncated like the long histories. The official test set and the SemEval-2027
  evaluation are next.
\end{itemize}

\section{Reproducibility}
Code: \todo{repository URL}. Setup and commands are in the README; every run writes its
configuration, git commit and upstream commit to \texttt{runs/<run id>/config.json}. The final
run is
\begin{center}\small
\texttt{python scripts/run\_baseline.py -{}-split dev-eval -{}-caps 2 3 4 -{}-retriever hybrid
-{}-fix-tool-schema}\\
\texttt{-{}-repair-tool-calls -{}-refresh-tools -{}-evidence-answer -{}-short-history -{}-policy-guard
-{}-tools-first}\\
\texttt{-{}-tool-hints -{}-llm-max-retries 8}
\end{center}
and the same with \texttt{-{}-caps 1 -{}-max-iterations 25} for Task~1 (an iteration limit as a cost
guard; the official runner sets none), each followed by \texttt{python scripts/evaluate.py -{}-run-id <run id>}. \todo{Final run ids.} The
ablation rows of Table~3 come from the runs listed in \texttt{results/v2\_and\_fixes\_dev\_eval.md}.

\textbf{Contributions.} \todo{Who did what.}

\end{document}
